% Generated by Listas/make_statements.py from pratica-lecture-06.md, item 5. Do not edit.
\textbf{Source.} Milgrom \& Roberts (1992), \emph{Economics, Organization and Management}, Ch. 8 Exercises, Mathematical Exercise 4, p.~286 {[}PDF 301{]}. Reproduced verbatim. Answer in the source: none in the book.



\begin{quote}

\begin{enumerate}

\def\labelenumi{\arabic{enumi}.}

\setcounter{enumi}{3}

\tightlist

\item

  In the discussion of the game of offering trust in the text, it was shown that if the game is being played only once, then trust will not be offered and, if it should be offered, it will not be honored. Suppose instead the two parties know they will interact exactly twice. Would trust be offered and honored in the second, final interaction? Why or why not? If the parties anticipate this and trust is offered in the first round, will it be honored? Why or why not? {[}Hint: Will honoring trust affect behavior at the last round?{]} Will trust then be offered? By this logic, what would happen if the parties knew they would interact exactly three times? Any arbitrary, but finite and known number of times? Do you find this to be a credible prediction about how people actually would behave? Why or why not?

\end{enumerate}

\end{quote}

